% Автоматаар үүсгэсэн — эх файл: lab03/README.md
\chapter{Лаб 3 — MQTT 5.0, Unified Namespace ба гурван зам}

\begingroup\scriptsize\begin{longtable}[]{@{}
  >{\raggedright\arraybackslash}p{(\columnwidth - 2\tabcolsep) * \real{0.5000}}
  >{\raggedright\arraybackslash}p{(\columnwidth - 2\tabcolsep) * \real{0.5000}}@{}}
\toprule
\endhead
\bottomrule
\endlastfoot
\textbf{7 хоног} & VII \\
\textbf{Хугацаа} & 4 цаг \\
\textbf{Суралцахуйн үр дүн} & ҮД1 (архитектур шинжлэх), ҮД4 (өгөгдлийн загвар зохиох), ҮД6 \\
\textbf{Үнэлгээ} & Бичгийн тайлан, 10 оноо \\
\textbf{Гол хэмжилт} & QoS × зам (loopback / LAN / bridge): p50, p95, \textbf{p99} \\
\end{longtable}\endgroup

\hypertarget{ux437ux43eux440ux438ux43bux433ux43e}{%
\section{1. Зорилго}\label{ux437ux43eux440ux438ux43bux433ux43e}}

Лаб 1-д гурван замын саатлыг \textbf{урьдчилан} нэг удаа хэмжсэн. Өнөөдөр түүнийг бүтэн матриц болгож, дээр нь протоколын шийдвэрүүдийг нэмнэ. Дөрвөн зүйлийг \textbf{тоогоор} тогтооно:

\begin{enumerate}
\def\labelenumi{\arabic{enumi}.}
\tightlist
\item
  MQTT 5.0-ийн зургаан боломж юуг шийддэг вэ --- ямар үнэтэй вэ?
\item
  \textbf{QoS × зам} матриц: loopback / 100 Mbit LAN / гүүр гурван замд QoS 0/1/2 тус бүр хэдэн мс вэ?
\item
  Тэсрэлт (burst) ба тогтвортой урсгалын саатал яагаад хэдэн арав дахин зөрдөг вэ?
\item
  Багийн төслийн \textbf{UNS схем} ямар байх вэ?
\end{enumerate}

Лабын гол бүтээгдэхүүн хоёр: \textbf{багийн UNS схем} ба \textbf{тоон үндэслэлтэй QoS сонголт}. Хоёуланг нь Лаб 4-өөс эцэс хүртэл ашиглана.

\begin{quote}
\textbf{p99 бол гол тоо, дундаж биш.} Pi 3B-гийн Ethernet нь USB 2.0 дээр сууж, зурвасаа USB-тэй хуваадаг. Тиймээс p50 сайхан харагдаж байхад p99 хэд дахин том гарч болно. Хэрэглэгчийн мэдэрдэг зүйл, SLA-д бичигддэг зүйл нь p99.
\end{quote}

\hypertarget{ux443ux440ux44cux434ux447ux438ux43bux441ux430ux43d-ux43dux4e9ux445ux446ux4e9ux43b}{%
\section{2. Урьдчилсан нөхцөл}\label{ux443ux440ux44cux434ux447ux438ux43bux441ux430ux43d-ux43dux4e9ux445ux446ux4e9ux43b}}

\begin{itemize}
\tightlist
\item
  Лаб 1, 2 дууссан, \texttt{lab02-done} тэмдэглэгээ тавигдсан
\item
  Бие даалт VI: MQTT 5.0 стандартын 3.3 (PUBLISH), 4.7 (Topic names) уншсан
\item
  \texttt{mosquitto\_pub} / \texttt{mosquitto\_sub} хоёр хост дээр суусан
\end{itemize}

\begin{Shaded}
\begin{Highlighting}[]
\CommentTok{\# [ЗК] үүлний давхарга}
\BuiltInTok{cd}\NormalTok{ \textasciitilde{}/cnc302/stack }\KeywordTok{\&\&} \FunctionTok{make}\NormalTok{ up }\KeywordTok{\&\&} \FunctionTok{make}\NormalTok{ health     }\CommentTok{\# бүх мөр 200}
\FunctionTok{make}\NormalTok{ ip                                          }\CommentTok{\# CLOUD\_HOST{-}ыг тэмдэглэ}
\CommentTok{\# [Pi] ирмэгийн давхарга}
\BuiltInTok{cd}\NormalTok{ \textasciitilde{}/cnc302/edge }\KeywordTok{\&\&} \FunctionTok{make}\NormalTok{ check                   }\CommentTok{\# бүгд OK}
\FunctionTok{make}\NormalTok{ up }\KeywordTok{\&\&} \FunctionTok{make}\NormalTok{ link                             }\CommentTok{\# bridge/state → 1}
\FunctionTok{make}\NormalTok{ mem                                         }\CommentTok{\# throttle 0x0 байх ЁСТОЙ}
\CommentTok{\# [Pi] Pi{-}гийн терминал БҮРД хувьсагчаа ачаална (доорх коммандууд үүнийг хэрэглэнэ)}
\BuiltInTok{set} \AttributeTok{{-}a} \KeywordTok{\&\&} \BuiltInTok{source}\NormalTok{ \textasciitilde{}/cnc302/edge/.env }\KeywordTok{\&\&} \BuiltInTok{set}\NormalTok{ +a}
\BuiltInTok{echo} \StringTok{"}\VariableTok{$CLOUD\_HOST}\StringTok{ / }\VariableTok{$SITE}\StringTok{/}\VariableTok{$AREA}\StringTok{/}\VariableTok{$LINE}\StringTok{/}\VariableTok{$DEVICE\_ID}\StringTok{"}
\end{Highlighting}
\end{Shaded}

\texttt{bridge/state} нь \texttt{1} биш бол \texttt{docs/troubleshooting.md} §1 рүү. Throttle \texttt{0x0} биш бол хөргөөд дахин эхэл --- өнөөдрийн бүх саатлын тоо тэр үед хүчингүй.

\hypertarget{ux43eux43dux43eux43bux44bux43d-ux441ux430ux43dux443ux443ux43bux433ux430}{%
\section{3. Онолын сануулга}\label{ux43eux43dux43eux43bux44bux43d-ux441ux430ux43dux443ux443ux43bux433ux430}}

\hypertarget{qos-ux431ux43eux43b-ux43dux430ux439ux434ux432ux430ux440ux442ux430ux439-ux431ux430ux439ux434ux43bux44bux43d-ux431ux430ux442ux430ux43bux433ux430ux430-ux431ux438ux448-ux434ux430ux43cux436ux443ux443ux43bux430ux43bux442ux44bux43d-ux433ux44dux440ux44dux44d}{%
\subsection{QoS бол найдвартай байдлын баталгаа биш, дамжуулалтын гэрээ}\label{qos-ux431ux43eux43b-ux43dux430ux439ux434ux432ux430ux440ux442ux430ux439-ux431ux430ux439ux434ux43bux44bux43d-ux431ux430ux442ux430ux43bux433ux430ux430-ux431ux438ux448-ux434ux430ux43cux436ux443ux443ux43bux430ux43bux442ux44bux43d-ux433ux44dux440ux44dux44d}}

\begingroup\scriptsize\begin{longtable}[]{@{}lllll@{}}
\toprule
QoS & Гэрээ & Гар барилт & Давхардах уу & Алдагдах уу \\
\midrule
\endhead
\bottomrule
\endlastfoot
0 & at most once & PUBLISH & үгүй & \textbf{тийм} \\
1 & at least once & PUBLISH → PUBACK & \textbf{тийм} & үгүй \\
2 & exactly once & PUBLISH → PUBREC → PUBREL → PUBCOMP & үгүй & үгүй \\
\end{longtable}\endgroup

QoS 2 нь дөрвөн алхамт тул хамгийн удаан. Гэхдээ энэ нь \textbf{зөвхөн дараагийн брокер хүртэлх} баталгаа: ирмэгийн mosquitto хүлээж авсан нь InfluxDB-д бичигдсэн гэсэн үг биш --- тэр гинжийг Лаб 5-д тасална.

\hypertarget{ux433ux443ux440ux432ux430ux43d-ux437ux430ux43c}{%
\subsection{Гурван зам}\label{ux433ux443ux440ux432ux430ux43d-ux437ux430ux43c}}

\begingroup\scriptsize\begin{longtable}[]{@{}
  >{\raggedright\arraybackslash}p{(\columnwidth - 4\tabcolsep) * \real{0.3333}}
  >{\raggedright\arraybackslash}p{(\columnwidth - 4\tabcolsep) * \real{0.3333}}
  >{\raggedright\arraybackslash}p{(\columnwidth - 4\tabcolsep) * \real{0.3333}}@{}}
\toprule
\begin{minipage}[b]{\linewidth}\raggedright
\texttt{-\/-path}
\end{minipage} & \begin{minipage}[b]{\linewidth}\raggedright
Юу дамжина
\end{minipage} & \begin{minipage}[b]{\linewidth}\raggedright
Юуг харуулна
\end{minipage} \\
\midrule
\endhead
\bottomrule
\endlastfoot
\texttt{loopback} & Pi доторх mosquitto, сүлжээгүй & брокер өөрөө хэдэн мс иддэг \\
\texttt{lan} & Pi → 100 Mbit уплинк → үүлний EMQX → буцаж Pi & сүлжээний бодит үнэ \\
\texttt{bridge} & Pi-гийн mosquitto, зэрэг гүүрээр үүл рүү урсгаж байхад & гүүрний дамжуулалт локал хүргэлтэд хэр саад болох \\
\end{longtable}\endgroup

\begin{quote}
⚠ \texttt{tools/qos\_latency.py} нь нийтлэгч ба захиалагчийг \textbf{нэг процесс дотор} ажиллуулдаг тул цагийн зөрүү (clock skew) байхгүй --- гэхдээ хоёр хостыг хамарсан төгсгөл-төгсгөлийн сааталыг хэмжиж \textbf{чадахгүй}. \texttt{-\/-path} бол \textbf{шошго}: хэмжилтийг ямар нөхцөлд хийснийг CSV-д тэмдэглэнэ. Гүүрээр дамжсан эсэхийг \textbf{тоолж} батална (Алхам 3В).
\end{quote}

\hypertarget{burst-ux431ux430-steady-state}{%
\subsection{Burst ба steady-state}\label{burst-ux431ux430-steady-state}}

500 мессежийг нэг дор шидвэл (\texttt{-\/-interval\ 0}) хэмжсэн "саатал" нь үнэндээ \textbf{дараалалд хүлээсэн хугацаа}. Тогтвортой урсгалд (\texttt{-\/-interval\ 0.01}) хэмжсэн нь \textbf{жинхэнэ сүлжээ + брокерын саатал}. Эхнийх нь хүчин чадлын хязгаарыг, хоёр дахь нь SLA-г тодорхойлно. Хоёрыг хольж тайлбарлах нь энэ лабораторийн хамгийн түгээмэл алдаа.

\hypertarget{unified-namespace-ux431ux43eux43b-ux442ux435ux445ux43dux43eux43bux43eux433ux438-ux431ux438ux448-ux433ux44dux440ux44dux44d}{%
\subsection{Unified Namespace бол технологи биш, гэрээ}\label{unified-namespace-ux431ux43eux43b-ux442ux435ux445ux43dux43eux43bux43eux433ux438-ux431ux438ux448-ux433ux44dux440ux44dux44d}}

"Аливаа өгөгдөл яг нэг л газар, урьдчилан таамаглах боломжтой нэрээр байрлана." Схемгүй бол 2 жилийн дараа брокер дээр \texttt{temp}, \texttt{Temperature}, \texttt{temp\_c}, \texttt{sensor/1/t} гэсэн дөрвөн нэр зэрэг оршино.

\begin{verbatim}
cnc302/<site>/<area>/<line>/<device>/<channel>
       shutis  mhts   lab    pi3b-01  telemetry
\end{verbatim}

Гүүр нь \texttt{cnc302/\textless{}SITE\textgreater{}/\#} сэдвийг \textbf{л} үүл рүү дамжуулна. Угтвар зөрвөл мессеж локал брокер дээр л үлдэнэ --- энэ курсын хамгийн түгээмэл ганц алдаа.

\hypertarget{ux430ux43bux445ux43cux443ux443ux434}{%
\section{4. Алхмууд}\label{ux430ux43bux445ux43cux443ux443ux434}}

\hypertarget{ux430ux43bux445ux430ux43c-1-ux441ux443ux443ux440ux44c-ux442ux4e9ux43bux4e9ux432-ux431ux430-ux445ux43eux451ux440-ux431ux440ux43eux43aux435ux440ux44bux43d-ux431ux430ux442ux430ux43bux433ux430ux430-15-ux43cux438ux43d-piux437ux43a}{%
\subsection{Алхам 1 --- Суурь төлөв ба хоёр брокерын баталгаа (15 мин) {[}Pi{]}{[}ЗК{]}}\label{ux430ux43bux445ux430ux43c-1-ux441ux443ux443ux440ux44c-ux442ux4e9ux43bux4e9ux432-ux431ux430-ux445ux43eux451ux440-ux431ux440ux43eux43aux435ux440ux44bux43d-ux431ux430ux442ux430ux43bux433ux430ux430-15-ux43cux438ux43d-piux437ux43a}}

\begin{Shaded}
\begin{Highlighting}[]
\CommentTok{\# [Pi] Pi дээр — юу ажиллаж байна, хэдий чинээ RAM үлдсэн}
\BuiltInTok{cd}\NormalTok{ \textasciitilde{}/cnc302 }\KeywordTok{\&\&} \ExtensionTok{./tools/measure\_stack.sh} \AttributeTok{{-}{-}role}\NormalTok{ edge}
\CommentTok{\# [ЗК] үүл дээр — гүүрээр юу ирж байгааг харна (энэ цонхыг НЭЭЛТТЭЙ орхи)}
\ExtensionTok{mosquitto\_sub} \AttributeTok{{-}h}\NormalTok{ localhost }\AttributeTok{{-}t} \StringTok{\textquotesingle{}cnc302/\#\textquotesingle{}} \AttributeTok{{-}v}
\end{Highlighting}
\end{Shaded}

{[}Pi{]} өөр терминалд агентыг асаа: \texttt{cd\ \textasciitilde{}/cnc302/edge\ \&\&\ make\ agent}. Компьютер дээрх цонхонд \texttt{telemetry}, \texttt{health}, \texttt{status} урсаж эхлэх ёстой.

\hypertarget{ux445ux4afux441ux43dux44dux433ux442-3.1-ux44dux445ux43bux44dux43bux438ux439ux43d-ux442ux4e9ux43bux4e9ux432}{%
\subsubsection{Хүснэгт 3.1 --- Эхлэлийн төлөв}\label{ux445ux4afux441ux43dux44dux433ux442-3.1-ux44dux445ux43bux44dux43bux438ux439ux43d-ux442ux4e9ux43bux4e9ux432}}

\begingroup\scriptsize\begin{longtable}[]{@{}lll@{}}
\toprule
Хэмжигдэхүүн & Утга & Тэмдэглэл \\
\midrule
\endhead
\bottomrule
\endlastfoot
Pi \texttt{MemAvailable} (MiB) & & Лаб 1-ийн Хүснэгт 1.1-тэй харьцуул \\
Pi температур (°C) / \texttt{get\_throttled} & / & \textbf{0x0} байх ёстой \\
\texttt{bridge/state} & & 1 байх ёстой \\
Үүл дээр харагдаж буй сувгууд & & telemetry / health / status / \ldots{} \\
\texttt{CLOUD\_HOST} & & \\
\texttt{SITE/\allowbreak{}AREA/\allowbreak{}LINE/\allowbreak{}DEVICE\_ID} & & \\
\end{longtable}\endgroup

\hypertarget{ux430ux43bux445ux430ux43c-2-mqtt-5.0-ux438ux439ux43d-ux437ux443ux440ux433ux430ux430ux43d-ux431ux43eux43bux43eux43cux436-ux445ux43eux451ux440-ux431ux440ux43eux43aux435ux440-ux434ux44dux44dux440-45-ux43cux438ux43d-ux437ux43api}{%
\subsection{Алхам 2 --- MQTT 5.0-ийн зургаан боломж, хоёр брокер дээр (45 мин) {[}ЗК{]}{[}Pi{]}}\label{ux430ux43bux445ux430ux43c-2-mqtt-5.0-ux438ux439ux43d-ux437ux443ux440ux433ux430ux430ux43d-ux431ux43eux43bux43eux43cux436-ux445ux43eux451ux440-ux431ux440ux43eux43aux435ux440-ux434ux44dux44dux440-45-ux43cux438ux43d-ux437ux43api}}

Эхлээд {[}ЗК{]} үүлний EMQX дээр дараалан ажиллуулж, гаралтыг \texttt{lab03/out/}-д хадгална:

\begin{Shaded}
\begin{Highlighting}[]
\BuiltInTok{cd}\NormalTok{ \textasciitilde{}/cnc302}
\ExtensionTok{python3}\NormalTok{ lab03/mqtt5\_features.py }\AttributeTok{{-}{-}host}\NormalTok{ localhost }\AttributeTok{{-}{-}expiry}\NormalTok{ 10 expiry}
\ExtensionTok{python3}\NormalTok{ lab03/mqtt5\_features.py }\AttributeTok{{-}{-}host}\NormalTok{ localhost }\AttributeTok{{-}{-}count}\NormalTok{ 1000 alias}
\ExtensionTok{python3}\NormalTok{ lab03/mqtt5\_features.py }\AttributeTok{{-}{-}host}\NormalTok{ localhost reqresp}
\ExtensionTok{python3}\NormalTok{ lab03/mqtt5\_features.py }\AttributeTok{{-}{-}host}\NormalTok{ localhost }\AttributeTok{{-}{-}subs}\NormalTok{ 3 }\AttributeTok{{-}{-}count}\NormalTok{ 300 shared}
\ExtensionTok{python3}\NormalTok{ lab03/mqtt5\_features.py }\AttributeTok{{-}{-}host}\NormalTok{ localhost }\AttributeTok{{-}{-}expiry}\NormalTok{ 60 session}
\ExtensionTok{python3}\NormalTok{ lab03/mqtt5\_features.py }\AttributeTok{{-}{-}host}\NormalTok{ localhost reason}
\end{Highlighting}
\end{Shaded}

Дараа нь \textbf{\texttt{expiry}, \texttt{shared}, \texttt{session} гурвыг {[}Pi{]} Pi дээр яг ижил коммандаар давтана} (\texttt{-\/-host\ localhost} → ирмэгийн mosquitto 2.0.20). Өөр програм хангамж, өөр зан төлөв --- зөрүүг Хүснэгт 3.2-т тэмдэглэ.

\hypertarget{ux445ux4afux441ux43dux44dux433ux442-3.2-mqtt-5.0-ux438ux439ux43d-ux431ux43eux43bux43eux43cux436ux443ux443ux434}{%
\subsubsection{Хүснэгт 3.2 --- MQTT 5.0-ийн боломжууд}\label{ux445ux4afux441ux43dux44dux433ux442-3.2-mqtt-5.0-ux438ux439ux43d-ux431ux43eux43bux43eux43cux436ux443ux443ux434}}

\begingroup\scriptsize\begin{longtable}[]{@{}
  >{\raggedright\arraybackslash}p{(\columnwidth - 10\tabcolsep) * \real{0.1667}}
  >{\raggedright\arraybackslash}p{(\columnwidth - 10\tabcolsep) * \real{0.1667}}
  >{\raggedright\arraybackslash}p{(\columnwidth - 10\tabcolsep) * \real{0.1667}}
  >{\raggedright\arraybackslash}p{(\columnwidth - 10\tabcolsep) * \real{0.1667}}
  >{\raggedright\arraybackslash}p{(\columnwidth - 10\tabcolsep) * \real{0.1667}}
  >{\raggedright\arraybackslash}p{(\columnwidth - 10\tabcolsep) * \real{0.1667}}@{}}
\toprule
\begin{minipage}[b]{\linewidth}\raggedright
Боломж
\end{minipage} & \begin{minipage}[b]{\linewidth}\raggedright
Ямар асуудлыг шийдэв
\end{minipage} & \begin{minipage}[b]{\linewidth}\raggedright
3.1.1-д хэрхэн шийддэг байсан
\end{minipage} & \begin{minipage}[b]{\linewidth}\raggedright
EMQX ({[}ЗК{]}) үр дүн
\end{minipage} & \begin{minipage}[b]{\linewidth}\raggedright
mosquitto ({[}Pi{]}) үр дүн
\end{minipage} & \begin{minipage}[b]{\linewidth}\raggedright
Бидний төсөлд
\end{minipage} \\
\midrule
\endhead
\bottomrule
\endlastfoot
Message Expiry & & & & & \\
Topic Alias & & & --- & --- & \\
Request/Response & & & & --- & \\
Shared Subscription & & & & & \\
Session Expiry & & & & & \\
Reason Code & & & & --- & \\
\end{longtable}\endgroup

\textbf{Alias-ийн хэмнэлтийг бодитоор хэмжих.} EMQX самбар (\texttt{:18083}) → \textbf{Monitoring → Metrics} → \texttt{Bytes\ Received}-ийг тэмдэглэ, доорхийг ажиллуул, дахин тэмдэглэ. Дараа нь \texttt{alias} туршилтыг давтаж гурав дахь удаа тэмдэглэ:

\begin{Shaded}
\begin{Highlighting}[]
\CommentTok{\# [ЗК] alias{-}гүй, урт сэдэв — 1000 мессеж}
\ExtensionTok{python3}\NormalTok{ tools/sim\_device.py }\AttributeTok{{-}{-}target}\NormalTok{ cloud }\AttributeTok{{-}{-}host}\NormalTok{ localhost }\AttributeTok{{-}{-}devices}\NormalTok{ 1 }\DataTypeTok{\textbackslash{}}
    \AttributeTok{{-}{-}count}\NormalTok{ 1000 }\AttributeTok{{-}{-}interval}\NormalTok{ 0 }\AttributeTok{{-}{-}site}\NormalTok{ shutis }\AttributeTok{{-}{-}area}\NormalTok{ campus{-}north{-}building{-}a}
\end{Highlighting}
\end{Shaded}

\hypertarget{ux445ux4afux441ux43dux44dux433ux442-3.3-topic-alias-ux438ux439ux43d-ux445ux44dux43cux43dux44dux43bux442}{%
\subsubsection{Хүснэгт 3.3 --- Topic Alias-ийн хэмнэлт}\label{ux445ux4afux441ux43dux44dux433ux442-3.3-topic-alias-ux438ux439ux43d-ux445ux44dux43cux43dux44dux43bux442}}

\begingroup\scriptsize\begin{longtable}[]{@{}
  >{\raggedright\arraybackslash}p{(\columnwidth - 8\tabcolsep) * \real{0.2000}}
  >{\raggedright\arraybackslash}p{(\columnwidth - 8\tabcolsep) * \real{0.2000}}
  >{\raggedright\arraybackslash}p{(\columnwidth - 8\tabcolsep) * \real{0.2000}}
  >{\raggedright\arraybackslash}p{(\columnwidth - 8\tabcolsep) * \real{0.2000}}
  >{\raggedright\arraybackslash}p{(\columnwidth - 8\tabcolsep) * \real{0.2000}}@{}}
\toprule
\begin{minipage}[b]{\linewidth}\raggedright
\end{minipage} & \begin{minipage}[b]{\linewidth}\raggedright
Сэдвийн урт (Б)
\end{minipage} & \begin{minipage}[b]{\linewidth}\raggedright
Мессеж
\end{minipage} & \begin{minipage}[b]{\linewidth}\raggedright
\texttt{Bytes\ Received} зөрүү
\end{minipage} & \begin{minipage}[b]{\linewidth}\raggedright
Нэг мессежид (Б)
\end{minipage} \\
\midrule
\endhead
\bottomrule
\endlastfoot
Alias-гүй & & 1000 & & \\
Alias-тай & & 1000 & & \\
\textbf{Хэмнэлт \%} & & & & \\
\end{longtable}\endgroup

\begin{quote}
Скриптийн тооцоолсон онолын хэмнэлт ба брокерын самбарын бодит зөрүү таарч байна уу? Зөрсөн бол шалтгааныг тайланд бич (санамж: TCP толгой, PUBACK, keepalive).
\end{quote}

\hypertarget{ux430ux43bux445ux430ux43c-3-qos-ux437ux430ux43c-ux43cux430ux442ux440ux438ux446-55-ux43cux438ux43d-pi}{%
\subsection{Алхам 3 --- QoS × зам матриц (55 мин) {[}Pi{]}}\label{ux430ux43bux445ux430ux43c-3-qos-ux437ux430ux43c-ux43cux430ux442ux440ux438ux446-55-ux43cux438ux43d-pi}}

\textbf{Энэ бол лабораторийн гол хэмжилт:} гурван зам × гурван QoS = 9 мөр, бүгд \textbf{Pi дээрээс} --- бид ирмэгийн төхөөрөмжийн туршлагыг хэмжиж байна. \textbf{В} мөрөнд сэдэв нь гүүрээр дамжих угтвартай тул мессеж бүр локал хүргэлтээс гадна зэрэг үүл рүү ч урсана.

\begin{Shaded}
\begin{Highlighting}[]
\CommentTok{\# [Pi] бүгдийг Pi дээр, дараалан}
\BuiltInTok{cd}\NormalTok{ \textasciitilde{}/cnc302}
\CommentTok{\# А. loopback — сүлжээ огт оролцохгүй}
\ExtensionTok{python3}\NormalTok{ tools/qos\_latency.py }\AttributeTok{{-}{-}path}\NormalTok{ loopback }\AttributeTok{{-}{-}host}\NormalTok{ localhost }\AttributeTok{{-}{-}port}\NormalTok{ 1883 }\DataTypeTok{\textbackslash{}}
    \AttributeTok{{-}{-}sweep} \AttributeTok{{-}{-}count}\NormalTok{ 500 }\AttributeTok{{-}{-}payload}\NormalTok{ 200 }\AttributeTok{{-}{-}interval}\NormalTok{ 0.01 }\AttributeTok{{-}{-}csv}\NormalTok{ lab03/out/qos{-}loopback.csv}
\CommentTok{\# Б. lan — Pi → 100 Mbit уплинк → үүлний EMQX → буцаж Pi}
\ExtensionTok{python3}\NormalTok{ tools/qos\_latency.py }\AttributeTok{{-}{-}path}\NormalTok{ lan }\AttributeTok{{-}{-}host} \StringTok{"}\VariableTok{$CLOUD\_HOST}\StringTok{"} \AttributeTok{{-}{-}port}\NormalTok{ 1883 }\DataTypeTok{\textbackslash{}}
    \AttributeTok{{-}{-}sweep} \AttributeTok{{-}{-}count}\NormalTok{ 500 }\AttributeTok{{-}{-}payload}\NormalTok{ 200 }\AttributeTok{{-}{-}interval}\NormalTok{ 0.01 }\AttributeTok{{-}{-}csv}\NormalTok{ lab03/out/qos{-}lan.csv}
\CommentTok{\# В. bridge — локал брокер, гүүр зэрэг дамжуулж байхад}
\ExtensionTok{python3}\NormalTok{ tools/qos\_latency.py }\AttributeTok{{-}{-}path}\NormalTok{ bridge }\AttributeTok{{-}{-}host}\NormalTok{ localhost }\AttributeTok{{-}{-}port}\NormalTok{ 1883 }\DataTypeTok{\textbackslash{}}
    \AttributeTok{{-}{-}topic} \StringTok{"cnc302/}\VariableTok{$SITE}\StringTok{/}\VariableTok{$AREA}\StringTok{/}\VariableTok{$LINE}\StringTok{/}\VariableTok{$DEVICE\_ID}\StringTok{/telemetry"} \DataTypeTok{\textbackslash{}}
    \AttributeTok{{-}{-}sweep} \AttributeTok{{-}{-}count}\NormalTok{ 500 }\AttributeTok{{-}{-}payload}\NormalTok{ 200 }\AttributeTok{{-}{-}interval}\NormalTok{ 0.01 }\AttributeTok{{-}{-}csv}\NormalTok{ lab03/out/qos{-}bridge.csv}
\end{Highlighting}
\end{Shaded}

Мессеж үнэхээр үүлэнд хүрч байгааг \textbf{зэрэг тоолж} батална:

\begin{Shaded}
\begin{Highlighting}[]
\CommentTok{\# [ЗК] үүл дээр, В хэмжилт эхлэхийн ӨМНӨ}
\ExtensionTok{mosquitto\_sub} \AttributeTok{{-}h}\NormalTok{ localhost }\AttributeTok{{-}t} \StringTok{"cnc302/}\VariableTok{$SITE}\StringTok{/\#"} \AttributeTok{{-}C}\NormalTok{ 500 }\KeywordTok{|} \FunctionTok{wc} \AttributeTok{{-}l}    \CommentTok{\# → 500}
\end{Highlighting}
\end{Shaded}

\begin{quote}
\textbf{Юуг хэмжиж, юуг хэмжээгүйгээ мэд.} В мөр нь мессеж үүлэнд хүрэх хугацааг биш, гүүр зэрэг ажиллаж байх үеийн \textbf{локал хүргэлтийг} хэмжинэ. А ба В-гийн зөрүү = гүүрний дамжуулалт локал брокерт хэр үнэтэй тусч байгаа нь. Төгсгөл-төгсгөлийн бүрэн бүтэн байдлыг Лаб 5-д InfluxDB дотор шалгана.
\end{quote}

\hypertarget{ux445ux4afux441ux43dux44dux433ux442-3.4-qos-ux437ux430ux43c-500-ux43cux435ux441ux441ux435ux436-200-ux431---interval-0.01}{%
\subsubsection{\texorpdfstring{Хүснэгт 3.4 --- QoS × зам (500 мессеж, 200 Б, \texttt{-\/-interval\ 0.01})}{Хүснэгт 3.4 --- QoS × зам (500 мессеж, 200 Б, -\/-interval 0.01)}}\label{ux445ux4afux441ux43dux44dux433ux442-3.4-qos-ux437ux430ux43c-500-ux43cux435ux441ux441ux435ux436-200-ux431---interval-0.01}}

\begingroup\scriptsize\begin{longtable}[]{@{}
  >{\raggedright\arraybackslash}p{(\columnwidth - 16\tabcolsep) * \real{0.1111}}
  >{\raggedright\arraybackslash}p{(\columnwidth - 16\tabcolsep) * \real{0.1111}}
  >{\raggedright\arraybackslash}p{(\columnwidth - 16\tabcolsep) * \real{0.1111}}
  >{\raggedright\arraybackslash}p{(\columnwidth - 16\tabcolsep) * \real{0.1111}}
  >{\raggedright\arraybackslash}p{(\columnwidth - 16\tabcolsep) * \real{0.1111}}
  >{\raggedright\arraybackslash}p{(\columnwidth - 16\tabcolsep) * \real{0.1111}}
  >{\raggedright\arraybackslash}p{(\columnwidth - 16\tabcolsep) * \real{0.1111}}
  >{\raggedright\arraybackslash}p{(\columnwidth - 16\tabcolsep) * \real{0.1111}}
  >{\raggedright\arraybackslash}p{(\columnwidth - 16\tabcolsep) * \real{0.1111}}@{}}
\toprule
\begin{minipage}[b]{\linewidth}\raggedright
Зам
\end{minipage} & \begin{minipage}[b]{\linewidth}\raggedright
QoS
\end{minipage} & \begin{minipage}[b]{\linewidth}\raggedright
p50 (мс)
\end{minipage} & \begin{minipage}[b]{\linewidth}\raggedright
p95 (мс)
\end{minipage} & \begin{minipage}[b]{\linewidth}\raggedright
\textbf{p99 (мс)}
\end{minipage} & \begin{minipage}[b]{\linewidth}\raggedright
max (мс)
\end{minipage} & \begin{minipage}[b]{\linewidth}\raggedright
Алдалт \%
\end{minipage} & \begin{minipage}[b]{\linewidth}\raggedright
Давхардал
\end{minipage} & \begin{minipage}[b]{\linewidth}\raggedright
Хүлээн авалт (мсж/с)
\end{minipage} \\
\midrule
\endhead
\bottomrule
\endlastfoot
loopback & 0 & & & & & & & \\
loopback & 1 & & & & & & & \\
loopback & 2 & & & & & & & \\
lan & 0 & & & & & & & \\
lan & 1 & & & & & & & \\
lan & 2 & & & & & & & \\
bridge & 0 & & & & & & & \\
bridge & 1 & & & & & & & \\
bridge & 2 & & & & & & & \\
\end{longtable}\endgroup

Матрицаас гурван харьцааг тооцож бич:

\begin{verbatim}
p99(lan, QoS1) / p99(loopback, QoS1)   = ______   ← сүлжээний үнэ
p99(bridge, QoS1) / p99(loopback,QoS1) = ______   ← гүүрний ачааллын үнэ
p99(QoS2) / p99(QoS1)  (нэг зам дээр)  = ______   ← гар барилтын үнэ
p99 / p50              (lan, QoS1)     = ______   ← СҮҮЛИЙН ӨРГӨН
\end{verbatim}

\begin{quote}
\textbf{Сүүлийн өргөн (p99/p50) нь Pi 3B-гийн гарын үсэг.} 100 Mbit Ethernet нь USB 2.0 контроллерын ард сууж байгаа тул диск, USB, сүлжээ нэг зурвасыг хуваана. p50 нь "бүх зүйл тайван байсан", p99 нь "USB завгүй байсан" агшнуудыг харуулна. Үйлдвэрийн хариу үйлдлийн баталгааг \textbf{сүүл} тодорхойлно.
\end{quote}

\hypertarget{ux430ux43bux445ux430ux43c-4-ux442ux44dux441ux440ux44dux43bux442-ux431ux430-ux442ux43eux433ux442ux432ux43eux440ux442ux43eux439-ux443ux440ux441ux433ux430ux43b-30-ux43cux438ux43d-pi}{%
\subsection{Алхам 4 --- Тэсрэлт ба тогтвортой урсгал (30 мин) {[}Pi{]}}\label{ux430ux43bux445ux430ux43c-4-ux442ux44dux441ux440ux44dux43bux442-ux431ux430-ux442ux43eux433ux442ux432ux43eux440ux442ux43eux439-ux443ux440ux441ux433ux430ux43b-30-ux43cux438ux43d-pi}}

Ижил зам (\texttt{lan}), зөвхөн илгээх хэмнэл ба ачаалал өөрчлөгдөнө.

\begin{Shaded}
\begin{Highlighting}[]
\CommentTok{\# [Pi] тэсрэлт: 500 мессежийг завсаргүй шиднэ}
\ExtensionTok{python3}\NormalTok{ tools/qos\_latency.py }\AttributeTok{{-}{-}path}\NormalTok{ lan }\AttributeTok{{-}{-}host} \StringTok{"}\VariableTok{$CLOUD\_HOST}\StringTok{"} \DataTypeTok{\textbackslash{}}
    \AttributeTok{{-}{-}sweep} \AttributeTok{{-}{-}count}\NormalTok{ 500 }\AttributeTok{{-}{-}payload}\NormalTok{ 200 }\AttributeTok{{-}{-}interval}\NormalTok{ 0 }\AttributeTok{{-}{-}csv}\NormalTok{ lab03/out/qos{-}burst.csv}
\CommentTok{\# [Pi] том ачаалал: 2000 байт, тогтвортой}
\ExtensionTok{python3}\NormalTok{ tools/qos\_latency.py }\AttributeTok{{-}{-}path}\NormalTok{ lan }\AttributeTok{{-}{-}host} \StringTok{"}\VariableTok{$CLOUD\_HOST}\StringTok{"} \DataTypeTok{\textbackslash{}}
    \AttributeTok{{-}{-}qos}\NormalTok{ 1 }\AttributeTok{{-}{-}count}\NormalTok{ 500 }\AttributeTok{{-}{-}payload}\NormalTok{ 2000 }\AttributeTok{{-}{-}interval}\NormalTok{ 0.01 }\AttributeTok{{-}{-}csv}\NormalTok{ lab03/out/qos{-}2000b.csv}
\CommentTok{\# [Pi] хуваалцсан захиалга: 3 хэрэглэгч}
\ExtensionTok{python3}\NormalTok{ tools/qos\_latency.py }\AttributeTok{{-}{-}path}\NormalTok{ lan }\AttributeTok{{-}{-}host} \StringTok{"}\VariableTok{$CLOUD\_HOST}\StringTok{"} \DataTypeTok{\textbackslash{}}
    \AttributeTok{{-}{-}qos}\NormalTok{ 1 }\AttributeTok{{-}{-}count}\NormalTok{ 500 }\AttributeTok{{-}{-}interval}\NormalTok{ 0.01 }\AttributeTok{{-}{-}shared}\NormalTok{ 3 }\AttributeTok{{-}{-}csv}\NormalTok{ lab03/out/qos{-}shared3.csv}
\end{Highlighting}
\end{Shaded}

\hypertarget{ux445ux4afux441ux43dux44dux433ux442-3.5-ux442ux44dux441ux440ux44dux43bux442-ux431ux430-ux442ux43eux433ux442ux432ux43eux440ux442ux43eux439-ux443ux440ux441ux433ux430ux43b-ux437ux430ux43c-lan}{%
\subsubsection{Хүснэгт 3.5 --- Тэсрэлт ба тогтвортой урсгал (зам = lan)}\label{ux445ux4afux441ux43dux44dux433ux442-3.5-ux442ux44dux441ux440ux44dux43bux442-ux431ux430-ux442ux43eux433ux442ux432ux43eux440ux442ux43eux439-ux443ux440ux441ux433ux430ux43b-ux437ux430ux43c-lan}}

\begingroup\scriptsize\begin{longtable}[]{@{}
  >{\raggedright\arraybackslash}p{(\columnwidth - 10\tabcolsep) * \real{0.1667}}
  >{\raggedright\arraybackslash}p{(\columnwidth - 10\tabcolsep) * \real{0.1667}}
  >{\raggedright\arraybackslash}p{(\columnwidth - 10\tabcolsep) * \real{0.1667}}
  >{\raggedright\arraybackslash}p{(\columnwidth - 10\tabcolsep) * \real{0.1667}}
  >{\raggedright\arraybackslash}p{(\columnwidth - 10\tabcolsep) * \real{0.1667}}
  >{\raggedright\arraybackslash}p{(\columnwidth - 10\tabcolsep) * \real{0.1667}}@{}}
\toprule
\begin{minipage}[b]{\linewidth}\raggedright
Горим
\end{minipage} & \begin{minipage}[b]{\linewidth}\raggedright
QoS
\end{minipage} & \begin{minipage}[b]{\linewidth}\raggedright
p50 (мс)
\end{minipage} & \begin{minipage}[b]{\linewidth}\raggedright
p99 (мс)
\end{minipage} & \begin{minipage}[b]{\linewidth}\raggedright
Илгээх хурд (мсж/с)
\end{minipage} & \begin{minipage}[b]{\linewidth}\raggedright
Хүлээн авалт (мсж/с)
\end{minipage} \\
\midrule
\endhead
\bottomrule
\endlastfoot
burst (\texttt{-\/-interval\ 0}) & 0 & & & & \\
burst & 1 & & & & \\
burst & 2 & & & & \\
steady (\texttt{-\/-interval\ 0.01}) & 0 & & & & \\
steady & 1 & & & & \\
steady & 2 & & & & \\
\end{longtable}\endgroup

burst / steady p50-ийн харьцаа (QoS 1): ule{1.5cm}{0.4pt} дахин. Шалтгаан: ule{1.5cm}{0.4pt}

\hypertarget{ux445ux4afux441ux43dux44dux433ux442-3.6-ux430ux447ux430ux430ux43bux430ux43b-ux431ux430-ux445ux443ux432ux430ux430ux43bux446ux441ux430ux43d-ux437ux430ux445ux438ux430ux43bux433ux44bux43d-ux43dux4e9ux43bux4e9ux4e9-qos-1-steady}{%
\subsubsection{Хүснэгт 3.6 --- Ачаалал ба хуваалцсан захиалгын нөлөө (QoS 1, steady)}\label{ux445ux4afux441ux43dux44dux433ux442-3.6-ux430ux447ux430ux430ux43bux430ux43b-ux431ux430-ux445ux443ux432ux430ux430ux43bux446ux441ux430ux43d-ux437ux430ux445ux438ux430ux43bux433ux44bux43d-ux43dux4e9ux43bux4e9ux4e9-qos-1-steady}}

\begingroup\scriptsize\begin{longtable}[]{@{}lllll@{}}
\toprule
Тохиргоо & p50 (мс) & p99 (мс) & Хүлээн авалт (мсж/с) & Тэмдэглэл \\
\midrule
\endhead
\bottomrule
\endlastfoot
200 Б, энгийн захиалга & & & & Хүснэгт 3.4-ээс \\
2000 Б, энгийн захиалга & & & & \\
200 Б, shared × 3 & & & & дараалал хадгалагдав уу? \\
\end{longtable}\endgroup

\begin{quote}
2000 байт нь 200 байтаас 10 дахин том. p99 нь 10 дахин өссөн үү? Үгүй бол хязгаарлагч нь зурвас биш, \textbf{нэг мессеж бүрийн тогтмол зардал} (round-trip, брокерын боловсруулалт) байна гэсэн үг. Энэ тоо Лаб 4-ийн уплинкийн тооцоонд шууд хэрэглэгдэнэ.
\end{quote}

\hypertarget{ux430ux43bux445ux430ux43c-5-unified-namespace-ux43cux43eux434-ux437ux4e9ux440ux447ux438ux43b-ux441ux445ux435ux43c-35-ux43cux438ux43d-piux437ux43a}{%
\subsection{Алхам 5 --- Unified Namespace: мод, зөрчил, схем (35 мин) {[}Pi{]}{[}ЗК{]}}\label{ux430ux43bux445ux430ux43c-5-unified-namespace-ux43cux43eux434-ux437ux4e9ux440ux447ux438ux43b-ux441ux445ux435ux43c-35-ux43cux438ux43d-piux437ux43a}}

\textbf{5.1 Одоогийн модыг харах.} Флотыг ирмэгийн брокер руу чиглүүлж, зэрэг сонсоно:

\begin{Shaded}
\begin{Highlighting}[]
\CommentTok{\# [ЗК] терминал 1 — виртуал флот Pi{-}гийн брокер руу}
\ExtensionTok{python3}\NormalTok{ tools/sim\_device.py }\AttributeTok{{-}{-}target}\NormalTok{ edge }\AttributeTok{{-}{-}host} \OperatorTok{\textless{}}\NormalTok{PI{-}IP}\OperatorTok{\textgreater{}}\NormalTok{ {-}{-}devices 8 }\AttributeTok{{-}{-}interval}\NormalTok{ 1 }\AttributeTok{{-}{-}lines}\NormalTok{ 3}
\CommentTok{\# [Pi] терминал 2 — модыг барьж авна}
\ExtensionTok{python3}\NormalTok{ lab03/uns\_tree.py }\AttributeTok{{-}{-}host}\NormalTok{ localhost }\AttributeTok{{-}{-}seconds}\NormalTok{ 30 }\DataTypeTok{\textbackslash{}}
    \AttributeTok{{-}{-}filter} \StringTok{\textquotesingle{}cnc302/\#\textquotesingle{}} \AttributeTok{{-}{-}json}\NormalTok{ lab03/out/uns{-}before.json}
\end{Highlighting}
\end{Shaded}

\textbf{5.2 Зөрчил үүсгэж, шалгагч барьж авахыг харах.} Шалгагч ажиллаж байх зуур {[}ЗК{]} дээрээс дараахийг илгээж, гурвуулангийн зөрчлийн тайланг тайландаа хуулна:

\begin{Shaded}
\begin{Highlighting}[]
\ExtensionTok{mosquitto\_pub} \AttributeTok{{-}h} \OperatorTok{\textless{}}\NormalTok{PI{-}IP}\OperatorTok{\textgreater{}}\NormalTok{ {-}t }\StringTok{\textquotesingle{}cnc302/UB/Campus/Line 1/DEV{-}1/temp\textquotesingle{}} \AttributeTok{{-}m} \StringTok{\textquotesingle{}\{"t":24\}\textquotesingle{}}
\ExtensionTok{mosquitto\_pub} \AttributeTok{{-}h} \OperatorTok{\textless{}}\NormalTok{PI{-}IP}\OperatorTok{\textgreater{}}\NormalTok{ {-}t }\StringTok{\textquotesingle{}sensors/room12/temperature\textquotesingle{}} \AttributeTok{{-}m} \StringTok{\textquotesingle{}24\textquotesingle{}}
\ExtensionTok{mosquitto\_pub} \AttributeTok{{-}h} \OperatorTok{\textless{}}\NormalTok{PI{-}IP}\OperatorTok{\textgreater{}}\NormalTok{ {-}t }\StringTok{\textquotesingle{}cnc302/ub/campus/line01/dev1/telemetry/raw\textquotesingle{}} \AttributeTok{{-}m} \StringTok{\textquotesingle{}\{\}\textquotesingle{}}
\end{Highlighting}
\end{Shaded}

\textbf{5.3 Хатуу горим ба гарах код.} \texttt{-\/-strict} нь CI-д ашиглагдана:

\begin{Shaded}
\begin{Highlighting}[]
\CommentTok{\# [Pi] Pi дээр, агент ажиллаж байх зуур}
\ExtensionTok{python3}\NormalTok{ lab03/uns\_tree.py }\AttributeTok{{-}{-}host}\NormalTok{ localhost }\AttributeTok{{-}{-}seconds}\NormalTok{ 20 }\AttributeTok{{-}{-}strict} \AttributeTok{{-}{-}prefix}\NormalTok{ cnc302}
\BuiltInTok{echo} \StringTok{"гарах код: }\VariableTok{$?}\StringTok{"}      \CommentTok{\# 0 = цэвэр, 1 = зөрчилтэй}
\end{Highlighting}
\end{Shaded}

\begin{quote}
\textbf{Хүлээгдэх ба сургамжтай үр дүн:} шалгагч нь \texttt{health}, \texttt{anomaly}, \texttt{config}, \texttt{bridge/state} сувгуудыг \textbf{зөрчил} гэж тэмдэглэнэ. Учир нь \texttt{uns\_tree.py}-ийн зөвшөөрөгдсөн сувгийн жагсаалт нь \texttt{telemetry\textbar{}status\textbar{}cmd\textbar{}event\textbar{}attributes}, мөн \texttt{bridge/state} нь 6 биш 7 хэсэгтэй. Энэ бол алдаа биш, \textbf{шийдвэр гаргах цэг}: та сувгийн жагсаалтыг өргөтгөх үү, эсвэл сувгуудаа схемд багтаах уу? Тайландаа сонголтоо үндэслэлтэй бич.
\end{quote}

\textbf{5.4 Багийн UNS схем.} \texttt{lab03/uns-schema.md} файлд бичнэ. Дор хаяж \textbf{10 бодит сэдвийн жишээ} оруулна:

\begin{itemize}
\tightlist
\item
  Таны төслийн шатлал (site / area / line / device) юу гэсэн үг вэ?
\item
  Ямар суваг хэрэгтэй вэ, тус бүр ямар чиглэлд урсах вэ?
\item
  Аль сэдэв \textbf{retained} байх вэ, яагаад? (Санамж: \texttt{status} бол LWT-тэй retained.)
\item
  Тушаал хаашаа очиж, хариу хаанаас ирэх вэ? Гүүрний \texttt{in} дүрэмтэй нийцэж байна уу?
\item
  10 жилийн дараа шинэ талбай нэмэхэд схем эвдрэх үү?
\end{itemize}

\hypertarget{ux445ux4afux441ux43dux44dux433ux442-3.7-ux431ux430ux433ux438ux439ux43d-uns-ux441ux445ux435ux43c}{%
\subsubsection{Хүснэгт 3.7 --- Багийн UNS схем}\label{ux445ux4afux441ux43dux44dux433ux442-3.7-ux431ux430ux433ux438ux439ux43d-uns-ux441ux445ux435ux43c}}

\begingroup\scriptsize\begin{longtable}[]{@{}
  >{\raggedright\arraybackslash}p{(\columnwidth - 10\tabcolsep) * \real{0.1667}}
  >{\raggedright\arraybackslash}p{(\columnwidth - 10\tabcolsep) * \real{0.1667}}
  >{\raggedright\arraybackslash}p{(\columnwidth - 10\tabcolsep) * \real{0.1667}}
  >{\raggedright\arraybackslash}p{(\columnwidth - 10\tabcolsep) * \real{0.1667}}
  >{\raggedright\arraybackslash}p{(\columnwidth - 10\tabcolsep) * \real{0.1667}}
  >{\raggedright\arraybackslash}p{(\columnwidth - 10\tabcolsep) * \real{0.1667}}@{}}
\toprule
\begin{minipage}[b]{\linewidth}\raggedright
Суваг
\end{minipage} & \begin{minipage}[b]{\linewidth}\raggedright
Чиглэл
\end{minipage} & \begin{minipage}[b]{\linewidth}\raggedright
QoS
\end{minipage} & \begin{minipage}[b]{\linewidth}\raggedright
Retained
\end{minipage} & \begin{minipage}[b]{\linewidth}\raggedright
Гүүрээр дамжих уу (\texttt{out}/\texttt{in}/үгүй)
\end{minipage} & \begin{minipage}[b]{\linewidth}\raggedright
Жишээ сэдэв
\end{minipage} \\
\midrule
\endhead
\bottomrule
\endlastfoot
telemetry & ирмэг → үүл & & & & \\
health & ирмэг → үүл & & & & \\
status & ирмэг → үүл & & & & \\
anomaly & ирмэг → үүл & & & & \\
bridge/state & ирмэг → үүл & & & & \\
cmd & үүл → ирмэг & & & & \\
config & үүл → ирмэг & & & & \\
ota/* & хоёр тал & & & & \\
\end{longtable}\endgroup

\hypertarget{ux430ux43bux445ux430ux43c-6-sparkplug-b-30-ux43cux438ux43d-ux437ux43a}{%
\subsection{Алхам 6 --- Sparkplug B (30 мин) {[}ЗК{]}}\label{ux430ux43bux445ux430ux43c-6-sparkplug-b-30-ux43cux438ux43d-ux437ux43a}}

⚠ \texttt{sparkplug\_lite.py} нь protobuf-гүй \textbf{хялбаршуулсан} хувилбар. Сэдвийн бүтэц, төлөвийн машин, \texttt{seq} нь стандартын дагуу; ачаалал нь JSON.

\textbf{6.1 Хоёр терминалд зэрэг ({[}ЗК{]} үүлний EMQX дээр):}

\begin{Shaded}
\begin{Highlighting}[]
\CommentTok{\# терминал 1 — ажиглагч; терминал 2 — ирмэгийн зангилаа}
\ExtensionTok{python3}\NormalTok{ lab03/sparkplug\_lite.py monitor }\AttributeTok{{-}{-}host}\NormalTok{ localhost }\AttributeTok{{-}{-}seconds}\NormalTok{ 90}
\ExtensionTok{python3}\NormalTok{ lab03/sparkplug\_lite.py node }\AttributeTok{{-}{-}host}\NormalTok{ localhost }\AttributeTok{{-}{-}devices}\NormalTok{ 3 }\AttributeTok{{-}{-}seconds}\NormalTok{ 60 }\AttributeTok{{-}{-}interval}\NormalTok{ 1}
\end{Highlighting}
\end{Shaded}

\texttt{NBIRTH\ →\ DBIRTH\ →\ DDATA\ …\ →\ NDEATH} дарааллыг ажигла. DDATA дотор метрикийн \textbf{нэр байхгүй, зөвхөн alias} байна --- ажиглагч нь BIRTH-ээс сурсан хүснэгтээрээ нэрийг сэргээнэ.

\textbf{6.2 Дарааллын тасалдал ба report-by-exception} (ажиглагч ажиллаж байх зуур):

\begin{Shaded}
\begin{Highlighting}[]
\ExtensionTok{python3}\NormalTok{ lab03/sparkplug\_lite.py node }\AttributeTok{{-}{-}host}\NormalTok{ localhost }\AttributeTok{{-}{-}devices}\NormalTok{ 2 }\AttributeTok{{-}{-}seconds}\NormalTok{ 30 }\AttributeTok{{-}{-}drop{-}seq}\NormalTok{ 5}
\ExtensionTok{python3}\NormalTok{ lab03/sparkplug\_lite.py node }\AttributeTok{{-}{-}host}\NormalTok{ localhost }\AttributeTok{{-}{-}devices}\NormalTok{ 3 }\AttributeTok{{-}{-}seconds}\NormalTok{ 30 }\AttributeTok{{-}{-}deadband}\NormalTok{ 0}
\ExtensionTok{python3}\NormalTok{ lab03/sparkplug\_lite.py node }\AttributeTok{{-}{-}host}\NormalTok{ localhost }\AttributeTok{{-}{-}devices}\NormalTok{ 3 }\AttributeTok{{-}{-}seconds}\NormalTok{ 30 }\AttributeTok{{-}{-}deadband}\NormalTok{ 0.5}
\end{Highlighting}
\end{Shaded}

\textbf{6.3 Гэнэтийн салалт (Last Will).} Зангилааг \textbf{Ctrl+C-гүйгээр} алж, NDEATH автоматаар гарахыг хүлээнэ:

\begin{Shaded}
\begin{Highlighting}[]
\ExtensionTok{python3}\NormalTok{ lab03/sparkplug\_lite.py node }\AttributeTok{{-}{-}host}\NormalTok{ localhost }\AttributeTok{{-}{-}devices}\NormalTok{ 2 }\AttributeTok{{-}{-}seconds}\NormalTok{ 300 }\KeywordTok{\&}
\FunctionTok{sleep}\NormalTok{ 15}
\BuiltInTok{kill} \AttributeTok{{-}9}\NormalTok{ \%1}
\end{Highlighting}
\end{Shaded}

\begin{quote}
Клиентийн keepalive нь кодод \textbf{30 секунд} бэхлэгдсэн тул NDEATH нь \texttt{kill\ -9}-ээс хойш \textbf{\textasciitilde30--45 секундын дараа} гарахыг хүлээнэ --- шууд биш. Энэ хоцролт нь "төхөөрөмж унасныг хэдэн секундын дараа мэдэх вэ" гэсэн бодит асуултын хариу юм.
\end{quote}

\textbf{6.4 Угтварын сургамж.} Sparkplug сэдэв нь \texttt{spBv1.0/…} угтвартай --- \texttt{cnc302/\textless{}SITE\textgreater{}/\#} \textbf{биш}. Тиймээс Pi дээр ажиллуулсан зангилааны мессежийг гүүр үүл рүү дамжуулахгүй. Батал:

\begin{Shaded}
\begin{Highlighting}[]
\CommentTok{\# [Pi] Pi дээр зангилаа ажиллуулж байхад, [ЗК] үүл дээр:}
\ExtensionTok{mosquitto\_sub} \AttributeTok{{-}h}\NormalTok{ localhost }\AttributeTok{{-}t} \StringTok{\textquotesingle{}spBv1.0/\#\textquotesingle{}} \AttributeTok{{-}v} \AttributeTok{{-}W}\NormalTok{ 15    }\CommentTok{\# хоосон байхыг хүлээнэ}
\end{Highlighting}
\end{Shaded}

\hypertarget{ux445ux4afux441ux43dux44dux433ux442-3.8-sparkplug-b-ux433ux438ux439ux43d-ux430ux436ux438ux433ux43bux430ux43bux442}{%
\subsubsection{Хүснэгт 3.8 --- Sparkplug B-гийн ажиглалт}\label{ux445ux4afux441ux43dux44dux433ux442-3.8-sparkplug-b-ux433ux438ux439ux43d-ux430ux436ux438ux433ux43bux430ux43bux442}}

\begingroup\scriptsize\begin{longtable}[]{@{}ll@{}}
\toprule
Асуулт & Хариулт \\
\midrule
\endhead
\bottomrule
\endlastfoot
NBIRTH-д хэдэн метрик зарлагдав & \\
DDATA-д метрикийн нэр дамжсан уу & \\
Дарааллын тасалдал илэрсэн үү (\texttt{-\/-drop-seq\ 5}) & \\
\texttt{kill\ -9}-ийн дараа NDEATH хэдэн секундын дараа гарав & \\
DDATA-гийн тоо: deadband 0 vs 0.5 & / \\
Deadband-ийн хэмнэлт \% & \\
Гүүр \texttt{spBv1.0/\#}-ийг дамжуулав уу & \\
\end{longtable}\endgroup

\hypertarget{ux430ux43bux445ux430ux43c-7-ux431ux430ux433ux438ux439ux43d-ux448ux438ux439ux434ux432ux44dux440-20-ux43cux438ux43d-ux437ux43a}{%
\subsection{Алхам 7 --- Багийн шийдвэр (20 мин) {[}ЗК{]}}\label{ux430ux43bux445ux430ux43c-7-ux431ux430ux433ux438ux439ux43d-ux448ux438ux439ux434ux432ux44dux440-20-ux43cux438ux43d-ux437ux43a}}

Энэ бол лабораторийн гол үр дүн: цаашид ямар тохиргоо хэрэглэхээ \textbf{тоон үндэслэлтэйгээр} шийднэ. Багана бүрд Хүснэгт 3.4--3.6-гийн \textbf{тодорхой тоо} иш татна ("хурдан учраас" гэсэн хариулт оноо авахгүй).

\hypertarget{ux445ux4afux441ux43dux44dux433ux442-3.9-ux431ux430ux433ux438ux439ux43d-ux43fux440ux43eux442ux43eux43aux43eux43bux44bux43d-ux448ux438ux439ux434ux432ux44dux440}{%
\subsubsection{Хүснэгт 3.9 --- Багийн протоколын шийдвэр}\label{ux445ux4afux441ux43dux44dux433ux442-3.9-ux431ux430ux433ux438ux439ux43d-ux43fux440ux43eux442ux43eux43aux43eux43bux44bux43d-ux448ux438ux439ux434ux432ux44dux440}}

\begingroup\scriptsize\begin{longtable}[]{@{}lll@{}}
\toprule
Шийдвэр & Сонголт & Тоон үндэслэл (аль хүснэгт, ямар тоо) \\
\midrule
\endhead
\bottomrule
\endlastfoot
Телеметрийн QoS & & \\
Тушаалын (cmd) QoS & & \\
Дохиоллын (anomaly) QoS & & \\
Session Expiry Interval & & \\
Message Expiry (телеметр) & & \\
Topic Alias хэрэглэх үү & & \\
Ачааллын дээд хэмжээ (Б) & & \\
Sparkplug B хэрэглэх үү & & \\
\end{longtable}\endgroup

\hypertarget{ux430ux43bux445ux430ux43c-8-git-commit-10-ux43cux438ux43d-ux437ux43a}{%
\subsection{Алхам 8 --- Git commit (10 мин) {[}ЗК{]}}\label{ux430ux43bux445ux430ux43c-8-git-commit-10-ux43cux438ux43d-ux437ux43a}}

\begin{Shaded}
\begin{Highlighting}[]
\BuiltInTok{cd}\NormalTok{ \textasciitilde{}/cnc302}
\FunctionTok{git}\NormalTok{ add lab03/}
\CommentTok{\# Хэмжилтийн нотолгоо: lab03/out/ нь .gitignore{-}д тул зориудаар албадан нэмнэ}
\FunctionTok{git}\NormalTok{ add }\AttributeTok{{-}f}\NormalTok{ lab03/out/qos{-}}\PreprocessorTok{*}\NormalTok{.csv lab03/out/uns{-}before.json}
\FunctionTok{git}\NormalTok{ status }\AttributeTok{{-}{-}short}
\FunctionTok{git}\NormalTok{ commit }\AttributeTok{{-}m} \StringTok{"Лаб 3: QoS × зам матриц, UNS схем, Sparkplug"}
\FunctionTok{git}\NormalTok{ tag lab03{-}done }\KeywordTok{\&\&} \FunctionTok{git}\NormalTok{ push }\KeywordTok{\&\&} \FunctionTok{git}\NormalTok{ push }\AttributeTok{{-}{-}tags}
\CommentTok{\# Нууц файл ороогүйг ЗААВАЛ шалга}
\FunctionTok{git}\NormalTok{ ls{-}files }\KeywordTok{|} \FunctionTok{grep} \AttributeTok{{-}E} \StringTok{\textquotesingle{}\textbackslash{}.env$|bridge\textbackslash{}.conf$|\textbackslash{}.key$|\textbackslash{}.crt$|devices\textbackslash{}.csv\textquotesingle{}}   \CommentTok{\# хоосон байх ЁСТОЙ}
\end{Highlighting}
\end{Shaded}

\hypertarget{ux445ux44fux43dux430ux43bux442ux44bux43d-ux430ux441ux443ux443ux43bux442ux443ux443ux434}{%
\section{5. Хяналтын асуултууд}\label{ux445ux44fux43dux430ux43bux442ux44bux43d-ux430ux441ux443ux443ux43bux442ux443ux443ux434}}

Хариулт бүр \textbf{хүснэгтээс авсан тоо} агуулсан байх ёстой.

\begin{enumerate}
\def\labelenumi{\arabic{enumi}.}
\item
  Хүснэгт 3.4-т \texttt{p99(lan)\ /\allowbreak{}\ p99(loopback)} харьцаа QoS 1 дээр хэд гарав? Тэр зөрүүний хэдэн хувь нь 100 Mbit уплинк, хэдэн хувь нь EMQX-ийн боловсруулалт вэ? Хэрхэн ялгаж тогтоох вэ?
\item
  Хүснэгт 3.4-т \texttt{p99/p50} харьцаа аль зам дээр хамгийн том байв? Хэрэв та SLA-д "95\% хүсэлт X мс дотор" гэж бичих бол X-ийг ямар тоогоор сонгох вэ, яагаад p50-ийг ашиглаж болохгүй вэ?
\item
  Хүснэгт 3.4-т QoS 2 нь QoS 1-ээс хэдэн дахин удаан байв? Тэр зөрүү нь зөвхөн гар барилтын нэмэлт хоёр алхмаас үүдэлтэй юу --- тоогоор шалга (нэг round-trip хэдэн мс байв?).
\item
  Хүснэгт 3.5-д burst ба steady горимын p50 хэдэн дахин зөрөв? Аль тоог нь брокерын гүйцэтгэл гэж нэрлэж \textbf{болохгүй} вэ, яагаад?
\item
  Хүснэгт 3.6-д ачааллыг 200 → 2000 байт болгоход p99 хэдэн хувиар өсөв? Хэрэв 10 дахин бага өссөн бол хязгаарлагч нь юу вэ? Энэ дүгнэлт Лаб 4-ийн уплинкийн тооцоонд хэрхэн нөлөөлөх вэ?
\item
  Хүснэгт 3.8-д \texttt{kill\ -9}-ээс NDEATH хүртэл хэдэн секунд өнгөрөв? Хэрэв таны төсөлд төхөөрөмж унасныг 5 секунд дотор мэдэх шаардлагатай бол ямар тохиргоог өөрчлөх вэ, тэр нь ямар үнэтэй вэ (Pi 3B-гийн 100 Mbit уплинк дээр 1000 төхөөрөмжөөр тооцоол)?
\item
  Хүснэгт 3.7-д \texttt{uns\_tree.py\ -\/-strict} таны сувгуудын хэдийг нь зөрчил гэж тэмдэглэв? Та схемээ өөрчлөх үү, шалгагчийг өөрчлөх үү? Багийн CI-д энэ шалгалтыг оруулбал ямар ашиг, ямар эрсдэл гарах вэ?
\end{enumerate}

\hypertarget{ux445ux4afux43bux44dux44dux43bux433ux44dux43d-ux4e9ux433ux4e9ux445-ux437ux4afux439ux43b}{%
\section{6. Хүлээлгэн өгөх зүйл}\label{ux445ux4afux43bux44dux44dux43bux433ux44dux43d-ux4e9ux433ux4e9ux445-ux437ux4afux439ux43b}}

\begingroup\scriptsize\begin{longtable}[]{@{}lll@{}}
\toprule
\# & Зүйл & Байрлал \\
\midrule
\endhead
\bottomrule
\endlastfoot
1 & Тайлан (Хүснэгт 3.1--3.9 бөглөсөн) & \texttt{lab03/report.md} → PDF \\
2 & Багийн UNS схем, 10 жишээ сэдэвтэй & \texttt{lab03/uns-schema.md} \\
3 & QoS хэмжилтийн CSV (5 файл) & \texttt{lab03/out/qos-*.csv} \\
4 & UNS модны JSON ба зөрчлийн гаралт & \texttt{lab03/\allowbreak{}out/\allowbreak{}uns-before.\allowbreak{}json} \\
5 & Sparkplug-ийн NBIRTH/DDATA/NDEATH гаралт & тайлангийн хавсралт \\
6 & Git tag & \texttt{lab03-done} \\
\end{longtable}\endgroup

\hypertarget{ux4afux43dux44dux43bux433ux44dux44dux43dux438ux439-ux448ux430ux43bux433ux443ux443ux440-10-ux43eux43dux43eux43e}{%
\section{7. Үнэлгээний шалгуур (10 оноо)}\label{ux4afux43dux44dux43bux433ux44dux44dux43dux438ux439-ux448ux430ux43bux433ux443ux443ux440-10-ux43eux43dux43eux43e}}

\begingroup\scriptsize\begin{longtable}[]{@{}
  >{\raggedright\arraybackslash}p{(\columnwidth - 2\tabcolsep) * \real{0.5000}}
  >{\raggedright\arraybackslash}p{(\columnwidth - 2\tabcolsep) * \real{0.5000}}@{}}
\toprule
\begin{minipage}[b]{\linewidth}\raggedright
Шалгуур
\end{minipage} & \begin{minipage}[b]{\linewidth}\raggedright
Оноо
\end{minipage} \\
\midrule
\endhead
\bottomrule
\endlastfoot
Хүснэгт 3.4 (QoS × зам, 9 мөр) бүрэн, throttle 0x0 үед хэмжигдсэн & 3 \\
Хүснэгт 3.5--3.6: burst/steady ялгааг \textbf{зөв тайлбарласан} & 2 \\
UNS схем зохиогдож, \texttt{-\/-strict} шалгагчаар туршигдсан & 2 \\
Sparkplug-ийн гурван ажиглалт (alias, seq gap, NDEATH хугацаа) & 1 \\
Хүснэгт 3.9 --- шийдвэр бүр тодорхой тоонд тулгуурласан & 1 \\
Git цэвэр, нууц файл ороогүй & 1 \\
\end{longtable}\endgroup

\begin{quote}
\textbf{Оноо хасагдах гол шалтгаан:} burst горимын саатлыг "брокерын саатал" гэж тайлбарлах, эсвэл p50-ийг гол тоо болгож бичих.
\end{quote}

\hypertarget{ux442ux4afux433ux44dux44dux43cux44dux43b-ux430ux43bux434ux430ux430}{%
\section{8. Түгээмэл алдаа}\label{ux442ux4afux433ux44dux44dux43cux44dux43b-ux430ux43bux434ux430ux430}}

\begingroup\scriptsize\begin{longtable}[]{@{}
  >{\raggedright\arraybackslash}p{(\columnwidth - 4\tabcolsep) * \real{0.3333}}
  >{\raggedright\arraybackslash}p{(\columnwidth - 4\tabcolsep) * \real{0.3333}}
  >{\raggedright\arraybackslash}p{(\columnwidth - 4\tabcolsep) * \real{0.3333}}@{}}
\toprule
\begin{minipage}[b]{\linewidth}\raggedright
Шинж
\end{minipage} & \begin{minipage}[b]{\linewidth}\raggedright
Шалтгаан
\end{minipage} & \begin{minipage}[b]{\linewidth}\raggedright
Шийдэл
\end{minipage} \\
\midrule
\endhead
\bottomrule
\endlastfoot
\texttt{bridge} мөрийн 500 мессеж үүлэнд ирэхгүй & сэдвийн угтвар \texttt{cnc302/\$SITE/} биш & \texttt{-\/-topic}-оо шалга; troubleshooting §2 \\
Бүх QoS дээр p50 ижил & \texttt{-\/-interval\ 0} --- бүгд дараалалд хүлээж байна & \texttt{-\/-interval\ 0.01} ашигла \\
QoS 2 дээр \texttt{АНХААР:\allowbreak{}\ …\ бүх\ мессеж\ ирсэнгүй} & \texttt{-\/-timeout} бага & \texttt{-\/-timeout\ 60} нэмнэ \\
Саатал 10 дахин хэлбэлзэнэ & throttling эсвэл Wi-Fi & \texttt{make\ mem} → 0x0 эсэх; кабельд шилж \\
\texttt{shared} туршилтад нэг захиалагч бүгдийг авна & \texttt{\$share} дэмжигдэхгүй & брокерын хувилбарыг шалга, хоёр брокер дээр давт \\
\texttt{expiry} "дараа нь ч ирлээ" гэнэ & брокер Message Expiry-г дэмжихгүй & Хүснэгт 3.2-т \textbf{ажиглалт болгон} бич \\
\texttt{uns\_tree.py} мессеж олохгүй & флот ажиллахгүй эсвэл \texttt{-\/-filter} буруу & \texttt{mosquitto\_sub\ -t\ \textquotesingle{}\#\textquotesingle{}\ -v} \\
\texttt{health}/\texttt{anomaly} зөрчил гэж гарна & шалгагчийн сувгийн жагсаалт хязгаартай & Алхам 5.3 --- шийдвэрээ тайланд бич \\
Sparkplug NDEATH гарахгүй & keepalive 30 сек дуусаагүй & 45 секунд хүлээ \\
\end{longtable}\endgroup
